% 02_related.tex — Related work / novelty. Drafted by item P2-E8-3.
% Voice: ../STYLE.md (Andreas Maier). Numbers: committed leaderboard / faithful_demo
% (\cite{leaderboardP2,faithfuldemoP2}). Figures: \ref{fig:faithfulness}, \ref{fig:taxonomy}.
% Honesty contract (plan.md §"What we claim about semantics"): we measure the semantic
% gap, we do not close it; all T3 labels are external and only causally verified.

\section{Related work}\label{sec:related}

Interpretability is studied in three communities that rarely cite one another, yet
all three ask the same question of a complex system: \emph{what, inside it, caused
this output?} We organise the prior work along those three traditions --- the
neuroscience battery, the attribution and explainable-AI toolkit, and mechanistic
interpretability --- and then turn to the literature that tries to \emph{validate}
such explanations. The recurring theme is the one our study removes: on the systems
these methods were built for, there is no ground truth, so faithfulness can be argued
but not measured. We close by positioning our delta against the few benchmarks that
\emph{do} supply ground truth, and against a different recent notion of what
``interpretable'' should mean.

\subsection*{Three traditions, one question}

The neuroscience tradition reads a system by recording it and perturbing it: tuning
curves, lesions, connectomics, dimensionality reduction, spike-train correlations,
field potentials, and Granger-style directed influence. Jonas and
Kording~\cite{jonas2017could} turned this battery on a known artifact --- a
microprocessor running classic games --- and showed that the methods produce rich,
publishable structure while missing how the chip actually computes. Their study is
the provocation behind ours, and we read it as a warning rather than a verdict: it
\emph{could not} score its own methods, because it reported what the analyses found,
not how far each fell from the known mechanism. We re-run the battery and attach that
missing number (Sec.~\ref{sec:results_A}).

The attribution and explainable-AI tradition asks, for a single decision, which inputs
mattered. Gradient and saliency maps~\cite{simonyan2014deep,sundararajan2017axiomatic,smilkov2017smoothgrad,springenberg2015striving},
class-activation methods~\cite{selvaraju2017gradcam,chattopadhyay2018gradcampp},
perturbation and occlusion methods~\cite{zeiler2014visualizing,fong2017interpretable,petsiuk2018rise},
and surrogate or game-theoretic attributions~\cite{ribeiro2016why,lundberg2017unified}
form the working toolkit. In the reinforcement-learning setting these methods have
been applied directly to Atari agents to visualise what a policy attends
to~\cite{greydanus2018visualizing,such2019atarizoo}. Yet a counterfactual analysis of
those very saliency maps found them \emph{exploratory, not
explanatory}~\cite{atrey2020exploratory}: the highlighted regions need not be the
causes. The doubt is well placed, but on a learned policy it can only be argued, never
settled, because the true cause is unknown. We settle it on the VCS, where the cause is
computable (Sec.~\ref{sec:results_B}).

Mechanistic interpretability, the youngest of the three, aims past attribution at the
\emph{circuit} --- the internal subgraph that implements a behaviour. Activation and
path patching, causal mediation, and causal scrubbing localise an effect by
intervening on internal activations~\cite{vig2020investigating,meng2022locating,wang2022interpretability,conmy2023towards};
distributed alignment search and interchange interventions test whether a hypothesised
variable is causally realised~\cite{geiger2021causal,geiger2023finding};
sparse dictionaries and probes try to read the system's features and concepts off its
state~\cite{cunningham2023sparse,alain2017understanding,hewitt2019designing,belrose2023eliciting,nanda2023attribution}.
These are the most causal tools in the field, and on a real artifact they are also the
ones we find most faithful (Sec.~\ref{sec:results_C}). The open question is what a
recovered circuit \emph{is}: we show, with the same oracle, that it is a correct
data-flow graph and not yet the computed function.

\subsection*{Faithfulness, and the trouble with measuring it}

Across all three traditions the agreed quality criterion is \emph{faithfulness} --- an
explanation is good if its claimed causes are the true causes. Jacovi and
Goldberg~\cite{jacovi2020towards} make the criterion precise and observe that it must
be defined operationally and evaluated with diagnostics matched to the explanation
type. The diagnostics, however, are necessarily indirect when the ground truth is
absent. Adebayo et al.~\cite{adebayo2018sanity} sharpened this into a now-standard
\emph{sanity check}: a saliency method should at least change when the model's weights
or the labels are randomised, and several popular methods do not, behaving like edge
detectors regardless of what the network learned. Such checks are powerful and
necessary. Still, they are \emph{negative} tests --- they can expose a method that is
clearly wrong, but they cannot certify that a method is right, because there is no
reference map to compare against. The field therefore optimises a proxy it cannot
directly score. Our oracle supplies the reference map and converts faithfulness from an
argued property into a measured one: on the discrete sprite-position regime, the causal
method \texttt{activation\_patching} reaches faithfulness $1.000$ while the popular
\texttt{vanilla\_saliency} collapses to $0.000$ on every one of the six core
games~\cite{faithfuldemoP2}, and the cross-tradition leaderboard plots all $31$ methods
against the same ground truth (Fig.~\ref{fig:faithfulness})~\cite{leaderboardP2}.

\subsection*{The semantic labels are external --- AtariARI and OCAtari}

Two public resources attach human-readable concepts to the Atari state. AtariARI
extracts labels for object positions, scores, and lives from commented
disassemblies and exposes them as a probing
benchmark~\cite{anand2019unsupervised}; OCAtari adds render-time offset
corrections and object detection so that the labelled coordinates align with what is
drawn~\cite{delfosse2023ocatari}. We use both, and we are deliberate about their
status. These labels are \emph{imported}, not discovered; our contribution is to
\emph{verify them causally} by intervention --- set the byte, re-render, confirm the
object moves to the predicted place on the bit-exact framebuffer --- which upgrades a
correlational label to a verified-causal one and auto-corrects its offset. This is the
hinge of our honesty contract: every concept-level (T3) score in this paper measures
\emph{alignment to a supplied label}, not the recovery of a meaning the method found on
its own. The distinction matters because a concept can be linearly decodable from the
state and still play no causal role, the ``present $\neq$ used'' trap that control
tasks were designed to catch~\cite{hewitt2019designing}; our intervention test is
strictly stronger than probing, and the resulting separation is one of the three
counted in Fig.~\ref{fig:taxonomy}.

\subsection*{Ground-truth benchmarks for interpretability --- and our delta}

Ours is not the first study to score interpretability against a known answer. The
existing benchmarks, however, build the answer \emph{into} the system. Yang and
Kim~\cite{yang2019benchmarking} construct images with a controlled, known relative
feature importance so that attribution maps can be graded against it, with explicit
false-positive metrics. Tracr~\cite{lindner2023tracr} compiles human-readable programs
into transformers whose circuits are therefore known by construction, and uses them as
a laboratory for mechanistic interpretability; InterpBench~\cite{gupta2024interpbench}
extends this to semi-synthetic transformers trained to realise a specified circuit,
giving a more realistic but still planted ground truth. These are valuable, and we
build on their logic. Our delta is twofold, and we state it without claiming an
unqualified ``first''. \emph{First}, our subject is a real, deployed, hand-authored,
complex artifact whose ground truth is \emph{independent of the artifact}: the VCS was
never designed to be interpretable, and we recover its causal structure from an
exhaustive intervention oracle on a bit-exact, differentiable
port~\cite{maier2025vcs}, the platform built on the established Atari evaluation
substrate~\cite{bellemare2013arcade} and the reference emulator~\cite{xitari}.
\emph{Second}, and more importantly, the planted-circuit benchmarks validate
\emph{faithfulness} --- did the method recover the circuit that was built in? --- and
nothing beyond it, because there the circuit \emph{is} the meaning by fiat. On a real
artifact the meaning is a separate object from the causal structure, so we can do
something none of these benchmarks can: use exact ground truth to \emph{separate
faithfulness from understanding}, and show that a method which perfectly recovers the
true causal wiring has not thereby recovered the semantics (Fig.~\ref{fig:taxonomy}).
We \emph{measure} that residual --- the semantic gap --- and are explicit that closing
it, by recovering a program's variables and design from scratch, is a different problem
and the subject of companion work.

\subsection*{A different ``interpretable''}

Finally, our use of the word \emph{interpretable} is one of several in circulation, and
we name the difference rather than blur it. We operationalise interpretability as
\emph{ground-truth recovery}: an explanation is interpretable to the degree that its
claimed structure matches the system's true causal structure, scored by the oracle.
Barbiero et al.~\cite{barbiero2025neural} instead define interpretability as
\emph{inference equivariance} --- a model is interpretable if a user can correctly
predict its outputs by simulating its reasoning --- and design networks to satisfy that
property. Their notion is a design target for new models; ours is an audit criterion
for an existing system whose answer is known. We cite their framework as a contrast,
not as the source of our definition. Read against Marr's levels of
analysis~\cite{marr1982vision} and Lazebnik's ``fix-the-radio''
parable~\cite{lazebnik2002radio}, both notions agree on the deeper point that drives
this paper: a faithful description of a system's wiring is not yet an account of the
algorithm it runs --- and on a fully known machine, that gap is no longer a matter of
opinion but a measured quantity.

% =====================================================================
% BIBTEX-NEEDED: full entries for the keys this section cites that are NOT
% yet declared in references.bib or in a sibling section's BIBTEX-NEEDED block.
% (SM to merge; do NOT edit references.bib here.)
%
% Already present in references.bib, reused as-is (NO entry below):
%   jonas2017could, simonyan2014deep (declared in 05_results_B BIBTEX-NEEDED),
%   sundararajan2017axiomatic, selvaraju2017gradcam, chattopadhyay2018gradcampp,
%   zeiler2014visualizing, fong2017interpretable, petsiuk2018rise, ribeiro2016why,
%   lundberg2017unified, greydanus2018visualizing, such2019atarizoo,
%   atrey2020exploratory, smilkov2017smoothgrad, springenberg2015striving,
%   bellemare2013arcade, xitari, delfosse2023ocatari.
%
% Already declared in a sibling section's BIBTEX-NEEDED block, reused as-is
% (NO entry below):
%   marr1982vision, adebayo2018sanity, anand2019unsupervised, hewitt2019designing,
%   vig2020investigating, meng2022locating, wang2022interpretability,
%   conmy2023towards, geiger2021causal, geiger2023finding, cunningham2023sparse,
%   alain2017understanding, belrose2023eliciting, nanda2023attribution,
%   maier2025vcs, leaderboardP2, faithfuldemoP2.
%
% NEW keys introduced by THIS section (SM: please add to references.bib):
%
% @inproceedings{jacovi2020towards,
%   author    = {Jacovi, Alon and Goldberg, Yoav},
%   title     = {Towards Faithfully Interpretable {NLP} Systems: How Should We
%                Define and Evaluate Faithfulness?},
%   booktitle = {Proceedings of the 58th Annual Meeting of the Association for
%                Computational Linguistics (ACL)},
%   year      = {2020},
%   pages     = {4198--4205},
%   publisher = {Association for Computational Linguistics},
%   doi       = {10.18653/v1/2020.acl-main.386}
% }
%
% @inproceedings{yang2019benchmarking,
%   author    = {Yang, Mengjiao and Kim, Been},
%   title     = {Benchmarking Attribution Methods with Relative Feature Importance},
%   booktitle = {NeurIPS Workshop / arXiv:1907.09701},
%   year      = {2019},
%   eprint        = {1907.09701},
%   archivePrefix = {arXiv},
%   primaryClass  = {cs.LG},
%   note      = {Introduces the BAM benchmark with known relative feature importance
%                and false-positive metrics for attribution methods.}
% }
%
% @inproceedings{lindner2023tracr,
%   author    = {Lindner, David and Kram{\'a}r, J{\'a}nos and Farquhar, Sebastian
%                and Rahtz, Matthew and McGrath, Thomas and Mikulik, Vladimir},
%   title     = {Tracr: Compiled Transformers as a Laboratory for Interpretability},
%   booktitle = {Advances in Neural Information Processing Systems (NeurIPS)},
%   volume    = {36},
%   year      = {2023},
%   eprint        = {2301.05062},
%   archivePrefix = {arXiv},
%   primaryClass  = {cs.LG}
% }
%
% @inproceedings{gupta2024interpbench,
%   author    = {Gupta, Rohan and Arcuschin, Iv{\'a}n and Kwa, Thomas and
%                Garriga-Alonso, Adri{\`a}},
%   title     = {{InterpBench}: Semi-Synthetic Transformers for Evaluating
%                Mechanistic Interpretability Techniques},
%   booktitle = {Advances in Neural Information Processing Systems (NeurIPS),
%                Datasets and Benchmarks Track},
%   volume    = {37},
%   year      = {2024},
%   eprint        = {2407.14494},
%   archivePrefix = {arXiv},
%   primaryClass  = {cs.LG}
% }
%
% @article{barbiero2025neural,
%   author    = {Barbiero, Pietro and Marra, Giuseppe and Ciravegna, Gabriele and
%                Debot, David and De Santis, Francesco and Giannini, Francesco and
%                Diligenti, Michelangelo and Lio, Pietro},
%   title     = {Neural Interpretable Reasoning},
%   journal   = {arXiv preprint arXiv:2502.11639},
%   year      = {2025},
%   eprint        = {2502.11639},
%   archivePrefix = {arXiv},
%   primaryClass  = {cs.LG}
% }
%
% @article{lazebnik2002radio,
%   author    = {Lazebnik, Yuri},
%   title     = {Can a Biologist Fix a Radio? --- Or, What I Learned While
%                Studying Apoptosis},
%   journal   = {Cancer Cell},
%   year      = {2002},
%   volume    = {2},
%   number    = {3},
%   pages     = {179--182},
%   doi       = {10.1016/S1535-6108(02)00133-2}
% }
% =====================================================================
